% appendix/oracle.tex -- paper-facing summary of ../harness/oracle/oracle_spec.md (W2/W25).
\section{Binding Oracle Details}\label{app:oracle}
For a case with base operation $b$ and fresh proof $\pi$ bound to $b$'s
canonical (\htm, \htu), the ground truth $\op(\pi)=\hat\exe(r_\pi)$ is measured
by sending $r_\pi$ directly to the executor with verification off and the
topology's override setting, including the base operation's canonical query or
body (e.g.\ \texttt{?op=freeze}). Records are joined on the proof's \jti; each
case mints a fresh proof, so each \jti belongs to one case.
Algorithm~\ref{alg:oracle} gives the hit predicate. Each hit records its
deviation kinds: \emph{method} (effective method $\neq$ base method),
\emph{route} (template differs), \emph{dispatch} (same template and method,
different operation) and \emph{resource} (same operation, different path
parameters, compared after percent-decoding). Severity is assigned from the
executed record: \textsf{S1} for a state-changing or privileged operation or a
different resource, \textsf{S2} for another read, \textsf{S3} for health or
public files. Hits are deduplicated by a root-cause signature (topology
factors, 1-minimal primitive set from delta debugging, class, deviation kind)
and confirmed by three reruns on freshly started topologies; a rerun denied for
clock skew (\texttt{iat\_skew}) is redone, not counted. The separate \GC
propagation predicate flags any client-chosen trust header that reaches the
executor; markers equal to the gateway-set subject are false positives and were
discounted.
